\documentclass[a4paper,12pt]{article}
\usepackage[utf8]{inputenc}
\usepackage{titlesec}
\usepackage[scale=0.8]{geometry}
\usepackage{graphicx}
\usepackage{xcolor,sectsty}

\definecolor{myred}{RGB}{128,0,32}
\subsectionfont{\color{myred}}
\sectionfont{\color{myred}}

\author{
        Ernandorena, Iván\\
        Legajo: E-1115/1\\
        \texttt{ivan.ernandorena@gmail.com}
}
\date{\today}
\title {
    \Huge\textsc{Trabajo Final - Agenda Personal\\}
    \large\textsc{Análisis de Lenguaje de Programación}
}

\begin{document}
  \pagenumbering{gobble}
  \thispagestyle{empty}
  \maketitle
  \begin{center}
  \includegraphics[scale=0.5]{photo}
  \end{center}

  \newpage
  \section*{Introducción}
    El trabajo final consiste en una agenda personal que además de llevar contactos como una agenda común, se
    pueden organizar reuniones, recordatorios, deudas y una lista de compras o cosas para hacer.
    \subsection*{Estructura de la agenda personal}
    La agenda consiste de:
    \begin{itemize}
      \item Un dueño: Este es quien creó la agenda y con este nombre se la guardara y se la cargara respectivamente.
      \item Contactos: Una lista de contactos de la forma \texttt{nombre-teléfono-dirección}.\\
      Por ejemplo: Martín 341-1235559 Pellegrini 200. También se pueden tener contactos solo de la forma \texttt{nombre-teléfono}.
      \item Recordatorios: También es una lista que puede tener \texttt{reuniones} y \texttt{recordatorios}.
      \item Deudas: Una lista de deudas de la forma \texttt{nombre-cantidad-razón}
      \item Compras o cosas para hacer: Por último se tiene una lista simple donde podemos escribir cualquier que se desee recordar.
    \end{itemize}

  \section*{Guia de usuario}
  \subsection*{Comandos}
    En una primera instancia tenemos los comandos para abrir, guardar, etc. Estos son los siguientes:
    \begin{itemize}
      \item \textbf{:? o :help} Imprime los demás comandos con sus respectivos argumentos.
      \item \textbf{:browse} Muestra el nombre de la agenda cargada actualmente, si hay alguna.
      \item \textbf{:load $\langle$nombre$\rangle$} Carga una agenda existente, utilizando el nombre con el que se la creo.
      \item \textbf{:display} Muestra toda la agenda cargada sin filtros.
      \item \textbf{:operations} Imprime la ayuda para las operaciones sobre la agenda.
      \item \textbf{:save} Guarda los cambios realizados.
      \item \textbf{:close} Cierra la agenda actualmente abierta. También guarda los cambios si se le indica.
      \item \textbf{:quit} Cierra la interfaz.
    \end{itemize}
    Todos los comandos pueden ser abreviados con la primera letra de cada uno. Por ejemplo, \texttt{:s} en lugar de \texttt{:save}.

    \subsection*{Operaciones}
    Luego, las operaciones sobre la agenda son la siguientes:
    \subsection*{Operaciones de creacion y agregado}
    \begin{itemize}
      \item \textbf{newSched $\langle$nombre$\rangle$} Crea una nueva agenda.
      
      \item \textbf{addContact $\langle$nombre$\rangle$ $\langle$teléfono$\rangle$ $\langle$dirección$\rangle$}
        Añade un nuevo contacto. Se puede añadir un contacto que tenga solo número de teléfono. Además tener en cuenta que el teléfono debe ser de la forma \texttt{caracteristica-número}.
      
      \item \textbf{addRemind $\langle$[R][M]$\rangle$ $\langle$día-mes-año$\rangle$ $\langle$hora:min$\rangle$ $\langle$descripción$\rangle$} Añade un nuevo recordatorio. Recordar añadir la R (reminder) o la M (meeting) al comienzo para poder diferenciarlos. Tambíen utilizar bien la notación para las fechas, por ejemplo 26-09-2017 o 26/09/2017; al igual que con la hora, utilizar los dos puntos.

      \item \textbf{addDebt $\langle$nombre$\rangle$ $\langle$\$dinero$\rangle$ $\langle$descripción$\rangle$}
      Añade una nueva deuda.

      \item \textbf{addGrocerie $\langle$descripción$\rangle$} Añade un nuevo elemento a la lista de compras.
    \end{itemize}
    \subsection*{Operaciones de eliminación}
    \begin{itemize}  
      \item \textbf{delContact $\langle$nombre$\rangle$} Elimina el contacto dado.
      \item \textbf{delRemind $\langle$[R][M]$\rangle$ $\langle$día-mes-año$\rangle$ $\langle$hora:min$\rangle$ $\langle$descripción$\rangle$} Elimina el recordatorio dado. Como se puede tener la misma descripción en varias fechas, o la misma fecha para distintas descripción está operación necesita todo el recordatorio.

      \item \textbf{delDebt $\langle$nombre$\rangle$ $\langle$\$dinero$\rangle$ $\langle$descripción$\rangle$} Elimina la deuda dada. También, como puede haber varias deudas de la misma persona se necesita toda la deuda.

      \item \textbf{delGrocerie $\langle$descripción$\rangle$} Elimina el elemento dado.
    \end{itemize}

    \subsection*{Operaciones de actualización}
    \begin{itemize}
      \item \textbf{updPhone $\langle$nombre$\rangle$ $\langle$nuevo teléfono$\rangle$} Actualiza el número de telefono de un contacto existente. Recordar como se ingresan los teléfonos en \textbf{addContact}

      \item \textbf{updAddress $\langle$nombre$\rangle$ $\langle$nueva dirección$\rangle$} Actualiza la dirección de un contacto existente.
    \end{itemize}

    \subsection*{Operaciones de busqueda y filtrado}
    \begin{itemize}
      \item \textbf{searchContact $\langle$nombre$\rangle$} Busca un contacto por su nombre.
      \item \textbf{allContacts} Lista todos los contactos.
      \item \textbf{allReminds} Lista todos los recordatorios.
      \item \textbf{allMeetings} Lista todas las reuniones.
      \item \textbf{allDebts} Lista todas las deudas.
      \item \textbf{allGroceries} Muestra la lista de compras.
      \item \textbf{interval $\langle$días$\rangle$} Lista todas las reuniones y recordatorios entre la fecha actual y $\langle$días$\rangle$
      \item \textbf{thisWeek} Lista todas las reuniones y recordatorios para la semana. Equivalente a realizar \texttt{interval 7}.
      \item \textbf{thisMonth} Idem. \textbf{thisWeek} pero para un mes.
      \item \textbf{debtsTo $\langle$nombre$\rangle$} Lista todas las deudas de $\langle$nombre$\rangle$.
      \item \textbf{debtsHigher $\langle$n$\rangle$} Muestra las deudas mayoras a $\langle$n$\rangle$.
    \end{itemize}

    \section*{Modulos del programa}
      \subsection*{Types.hs}
       Aquí se definieron todas las principales estructuras utilizadas, asi como también se redefinieron algunos tipos con la sentencia \texttt{type} para que luego el código sea más fácil de leer.\\
       Además de los tipos y las estructuras para la agenda, se tiene la estructura del estado que lleva el programa en todo momento. El estado consta de el último archivo cargado con el comando \texttt{:load}
       y de una instancia de la agenda cargada, si es que hay una. \\ 
       Para tener la posibilidad de no tener una agenda cargada se definió el siguiente tipo de dato:
       \begin{verbatim}
       data LoadSched = Null | LS Schedule
       \end{verbatim}
       donde \texttt{Schedule} es la estructura de la agenda definida en el mismo modulo y explicada más arriba.

       \subsection*{AST.hs}
       En este módulo se encuentra el árbol de sintaxis abstracta de nuestro lenguaje sobre la agenda.
       Simplemente estan los tipos que representan los comandos y las operaciones explicadas en la sección
       anterior. Cuando se lee, por ejemplo, la operación \texttt{newSched Juan} el analizador o \textit{parser}
       devuelve algo de la forma: \begin{verbatim} NewSched Juan :: ScheduleComm \end{verbatim} Idem. con las demás operaciones y comandos, con la excepción de que los comandos están separados en el tipo \textit{InterpreterComm}.

       \subsection*{}



\end{document}